% Reusable CV example. All people, organisations and achievements are fictional.
% Keep this file next to cv.cls. Replace the sample text and placeholder URLs.
% Build: latexmk -pdf cv-example.tex
% Use [letterpaper] for US Letter; A4 is the default.
\documentclass{cv}

% APPEARANCE
% All settings are optional. Sizes accept dimensions such as 10pt.
% Font choices: computer-modern, latin-modern, times, palatino, helvetica.
% body-size also controls heading, subentry, date, location, contact, bio and footer sizes
% unless those sizes are explicitly set below. Line spacing is 1.2 x size.
\cvsetup{
    font          = computer-modern,
    body-size     = 10pt,
    name-size     = 17.28pt,
    section-size  = 12pt,
    link-color    = blue,
    location-color = black!55
}
% Further optional settings (shown with their defaults):
% \cvsetup{
%     heading-size       = 10pt,
%     subentry-size      = 10pt,
%     date-size          = 10pt,
%     location-size      = 10pt,
%     contact-size       = 10pt,
%     bio-size           = 10pt,
%     footer-size        = 10pt,
%     location-icon-size = 7pt,
%     contact-icon-color = black!40,
%     organisation-first = true,
%     show-dates         = true,
%     show-location-icons = true,
%     show-contact-icons = true,
%     show-page-numbers  = true,
%     margin             = 2cm,
%     heading-gap        = 4pt,
%     bullet-gap         = 2pt,
%     entry-gap          = 7pt,
%     section-gap        = 10pt,
%     header-gap         = 6pt,
%     bio-gap            = 10pt,
%     date-gap           = 8pt,
%     location-gap       = 10pt,
%     subentry-indent    = 10pt,
%     row-inset          = 6pt,
%     date-width         = auto
% }
% date-width=auto measures a date range in the selected date font.
% For longer date labels, supply a width such as date-width=3cm.
% show-dates=false removes the entire date column and reclaims its width.
% One entry without dates still aligns with the other entries.
% Colours accept any xcolor expression, e.g. blue!70!black.
% The centred footer reads "page 1 of 1" (also shown on one-page CVs).
% latexmk resolves the total automatically; with pdflatex, compile twice.

% HEADER
% Every field is optional, including name. Omit fields or set them to {}.
% If no fields are provided, \makecvheader produces no space or separators.
\cvheader{
    name           = {Alex Example},
    title          = {Research Engineer},
    email          = {alex@example.com},
    website        = {https://example.com},
    github         = {https://github.com/your-name},
    linkedin       = {https://www.linkedin.com/in/your-name/},
    google-scholar = {https://scholar.google.com/citations?user=YOUR_ID},
    % phone        = {+41 00 000 00 00},
    % location     = {Zürich, Switzerland},
    bio            = {Research engineer working on information retrieval and machine learning.
                      Develops evaluation methods and turns research prototypes into reliable software.}
}
% Delete any unwanted field above; uncomment phone or location if useful.
% Keep the contact row short for the cleanest layout.
% Display labels can differ from their destinations; defaults are shown below.
% \cvheader{
%     email-label    = {Email},
%     website-label  = {Website},
%     github-label   = {GitHub},
%     linkedin-label = {LinkedIn},
%     google-scholar-label = {Google Scholar}
% }
% For example, website-label={Portfolio} changes only the visible link text.
% The optional bio is left-aligned below the contacts and has no section title.

\begin{document}
\makecvheader

% ENTRIES
% All five entry fields are optional and reset for every entry.
% Missing organisation or role: no middle dot.
% Missing location: no pin; the title uses the full available width.
% Missing URL: plain heading text. URL links the organisation, or the role
% if organisation is absent. Organisation names are automatically bold.
% Braces protect commas in field values; use \& for an ampersand.
% Entries stay together on a page; keep each entry shorter than one page.
\cvsection{Professional Experience}

% Complete entry: linked organisation, role, dates, location, detail and bullets.
\begin{cventry}[
    organisation = {Example Research},
    role         = {Research Engineer},
    dates        = {09/2024--present},
    location     = {Zürich, Switzerland},
    url          = {https://example.com}
]
    \cvdetail{Project: \cvlink{https://example.com/search}{\textit{Evaluating Search over Company Knowledge}}}
    \begin{cvbullets}
        \item Built evaluation frameworks for retrieval and document parsing; deployed the resulting features in production.
        \item Developed a benchmark spanning documentation, support conversations and source code.
    \end{cvbullets}
\end{cventry}

% Omit url for an unlinked heading; omit location to give the title more space.
\begin{cventry}[
    organisation = {Example University},
    role         = {Teaching Assistant},
    dates        = {09/2022--06/2024}
]
    \begin{cvbullets}
        \item Prepared and taught weekly mathematics exercise classes for 30 undergraduate students.
    \end{cvbullets}
\end{cventry}

\cvsection{Education}

\begin{cventry}[
    organisation = {Example University},
    role         = {MSc Computer Science},
    dates        = {09/2022--06/2024},
    location     = {Zürich, Switzerland},
    url          = {https://example.com/university}
]
    \cvdetail{GPA: 5.7/6.0; credits: 120/120 ECTS; focus: Machine Learning}
    \cvdetail{Thesis: \cvlink{https://example.com/thesis}{\textit{Learning to Retrieve Evidence from Technical Documents}}}
    \cvdetail{Supervisor: Prof. Jordan Example (Information Retrieval Group)}
    % Subentries take the same five fields, but have no body.
    % Titles indent under the degree; their dates stay in the main date column.
    \cvsubentry[
        organisation = {Partner University},
        role         = {Exchange Semester},
        dates        = {01/2023--06/2023},
        location     = {Copenhagen, Denmark}
    ]
    \cvsubentry[
        organisation = {Example Institute},
        role         = {Summer School},
        dates        = {08/2023},
        location     = {London, UK}
    ]
\end{cventry}

% A heading-only entry needs no body commands.
\begin{cventry}[
    organisation = {Example University},
    role         = {BSc Mathematics},
    dates        = {09/2019--06/2022}
]
\end{cventry}

% Section names are unrestricted: add, remove or rename them to suit your CV.
\cvsection{Selected Projects \& Publications}

% Organisation-only entry: bold project name, with no stray middle dot.
\begin{cventry}[
    organisation = {Open Search Toolkit},
    dates        = {2024--present},
    url          = {https://example.com/toolkit}
]
    \begin{cvbullets}
        \item Created an open-source toolkit for comparing retrieval methods on reproducible datasets.
    \end{cvbullets}
\end{cventry}

% A role-only entry is useful for a publication title. Its URL links that title.
\begin{cventry}[
    role  = {Evaluating Retrieval on Heterogeneous Collections},
    dates = {2024},
    url   = {https://example.com/paper}
]
    \cvdetail{\textbf{Alex Example}, Jordan Example; Example Research Workshop.}
    \cvdetail{\cvlink{https://example.com/paper}{Paper}; \cvlink{https://example.com/code}{Code}; \cvlink{https://example.com/data}{Dataset}}
\end{cventry}

\cvsection{Additional Experience}

% No dates: the entry still aligns with the content column above.
\begin{cventry}[role={Volunteer Programming Mentor}]
    \cvdetail{Led introductory Python workshops for local secondary-school students.}
\end{cventry}

% No heading fields: use a body-only entry for skills, languages or interests.
% Remove this entire section if these details are better shown in your entries.
\cvsection{Skills \& Languages}
\begin{cventry}
    \cvdetail{\textbf{Technical:} Python, SQL, PyTorch, Git, Linux}
    \cvdetail{\textbf{Languages:} English (fluent); German (native); French (conversational)}
\end{cventry}

% LAYOUT RECIPES (set these in \cvsetup in the preamble, not here)
% Role before organisation: organisation-first=false (organisation stays bold).
% No date column: show-dates=false. No icons: show-contact-icons=false,
% show-location-icons=false. No footer: show-page-numbers=false.
% Compact style: body-size=9pt, heading-gap=3pt, bullet-gap=1pt, entry-gap=5pt.
% Sans-serif style: font=helvetica. Muted links: link-color=blue!50!black.
%
% Longer CVs flow onto additional pages automatically, with page X of Y footers.
% Each cventry (including its subentries) stays together; it must fit on one page.
% To choose a page break yourself, place \newpage between complete entries.

\end{document}
